% !PRINTED 600
\weberpage{600}
If now we set

\[
\mathfrak{p}_{r,s} = \mathfrak{p}_r\,|\,\varphi_{r,s}
\]

then by \S.\,160, (6),

\weq{6}{\mathfrak{p}_r' = (\mathfrak{p}_{r,1}\,\mathfrak{p}_{r,2}\cdots\mathfrak{p}_{r,e_r})^{g_r},}

and the substitution of (6) in (3) yields, through comparison
with \S.\,160, (6),

\weq{7}{g = a_r\,g_r, \quad e_1 + e_2 + \cdots + e_{e'} = e,}

whereby the exponents $a_r$ are defined.

The partial norm of $\mathfrak{p}_{r,s}$, taken with respect to the field $\Omega'$,
is by \S.\,160, (3), (7):

\weq{8}{\mathfrak{N}_{\Omega'}(\mathfrak{p}_{r,s}) = \mathfrak{p}_r'^{\,f_r},}

if $f_r$ is defined by the equation

\weq{9}{n' = e_r\,f_r\,g_r, \quad (h_r = f_r\,g_r).}

We now wish to decompose the group $\Phi$ into subsidiary groups
with respect to $\Phi'$, and, if $\vartheta_1, \vartheta_2, \ldots, \vartheta_{m'}$ are suitably
chosen substitutions from $\Phi$, we set:

\weq{10}{\Phi = \Phi'\vartheta_1 + \Phi'\vartheta_2 + \cdots + \Phi'\vartheta_{m'},}

so that, by the substitutions $\vartheta_1, \vartheta_2, \ldots, \vartheta_{m'}$, the field
$\Omega'$ passes over into $m'$ conjugate (equal or different) fields

\[
\Omega_1', \ \Omega_2', \ \ldots, \ \Omega_{m'}'.
\]

The field $\Omega_t'$ then belongs to the group $\vartheta_t^{-1}\Phi'\vartheta_t$.

If now, by $\vartheta_t$, the ideals $\mathfrak{p}_r'$ and $\mathfrak{p}_{r,s}$ pass over into $\mathfrak{p}_{r,t}'$ and
$\mathfrak{p}_{r,s,t}$, then by \S.\,160, (6)

\weq{11}{\mathfrak{p}_{r,t}' = (\mathfrak{p}_{r,1,t}\,\mathfrak{p}_{r,2,t}\cdots\mathfrak{p}_{r,e_r,t})^{g_r},}

and the product of all these ideals for $t = 1, 2, \ldots, m'$ is the
partial norm of $\mathfrak{p}_r'$ taken in the field $\Omega'$ with respect to the
field $R$, which we denote by $\mathfrak{N}_R'(\mathfrak{p}_r')$. It must be a power
of $\mathfrak{P}$, and we set

\weq{12}{\mathfrak{N}_R'(\mathfrak{p}_r') = \mathfrak{P}^{\,f_r'}.}

To find $f_r'$, we need only decompose $\mathfrak{P}$ in (12), according to \S.\,160, (6),
into its prime factors $\mathfrak{p}$, whereby there result $e\,g\,f_r'$ prime factors
$\mathfrak{p}$ in $\mathfrak{P}^{\,f_r'}$. If, on the other hand, we form the product
of the $m'$ factors (11), we obtain $g_r\,e_r\,m'$ such prime factors.
Hence

\[
e\,g\,f_r' = m'\,g_r\,e_r,
\]

and therefore, by (1), (9) and \S.\,160, (7):

\weq{13}{f = f_r\,f_r'.}

\begin{center}
\rule{0.3\linewidth}{0.4pt}
\end{center}
